\documentclass[11pt,letterpaper]{article}

\usepackage[margin=2.2cm]{geometry}
\usepackage{amsmath}
\usepackage{amssymb}
\usepackage{array}
\usepackage{booktabs}
\usepackage{graphicx}
\usepackage{pdflscape}

\title{Memoria de calculo alternativa\\Proyecto IE0413}
\author{Daniel Gerardo Madriz Chavarr\'ia C24478\\
Fabricio Torres Jim\'enez C27803\\
Kenneth Daniel Vizca\'ino Jim\'enez C09331}
\date{I ciclo 2026}

\begin{document}

\maketitle

\section*{Nota}

Esta version evita el uso de un flip-flop externo. Para obtener la senal a
\(f_0/2\), se usa un segundo \texttt{NE555} en modo monoestable como divisor de
frecuencia. La justificacion del filtro se escribe usando conceptos vistos en
la materia: diagramas de Bode, frecuencia normalizada, cascada de etapas,
coeficientes \(a,b\), factor \(Q\), banda de paso y banda de rechazo.

En esta version el \texttt{NE555} se usa para la temporizacion principal:
oscilador cuadrado y monoestable divisor. Las etapas analogicas posteriores se
implementan con \texttt{LF353}: inversion o buffer de senales, cambio de nivel,
integrador y filtro activo. Esto mantiene el circuito dentro de bloques vistos
en el curso sin introducir un flip-flop externo, una fuente ideal de corriente
ni herramientas fuera de la materia vista.

\begin{landscape}
\section*{Esquematico general}
\begin{center}
\includegraphics[width=\linewidth,height=0.82\textheight,keepaspectratio]{figuras/esquematico_sin_cables_en_aire_ne555.png}
\end{center}
\end{landscape}

\section{Datos del grupo}

Los primeros apellidos se ordenan alfabeticamente:

\[
\text{Madriz}, \qquad \text{Torres}, \qquad \text{Vizca\'ino}.
\]

\[
X=2, \qquad Y=2, \qquad Z=9.
\]

La frecuencia central pedida es:

\[
f_0=
\left(
X+\frac{Y}{10}+\frac{Z}{100}
\right)1000~\text{Hz}.
\]

\[
f_0=
\left(
2+\frac{2}{10}+\frac{9}{100}
\right)1000.
\]

\[
f_0=(2+0.2+0.09)1000.
\]

\[
\boxed{f_0=2290~\text{Hz}=2.29~\text{kHz}}.
\]

Las frecuencias requeridas son:

\[
f_{saw}=2f_0=4580~\text{Hz}.
\]

\[
f_{tri}=\frac{f_0}{2}=1145~\text{Hz}.
\]

\section{Onda cuadrada con NE555}

El \texttt{NE555} se conecta como astable tipo Schmitt. Las entradas
\texttt{TRIG} y \texttt{THRES} se conectan al capacitor \(C_{osc}\). La salida
\texttt{OUT} carga y descarga el capacitor por medio de \(R_{osc}\).

Los niveles internos son:

\[
V_L=\frac{1}{3}V_{CC},
\qquad
V_H=\frac{2}{3}V_{CC}.
\]

Durante la carga:

\[
v_C(t)=V_{CC}+
\left(
\frac{1}{3}V_{CC}-V_{CC}
\right)e^{-t/(RC)}.
\]

\[
v_C(t)=V_{CC}-\frac{2}{3}V_{CC}e^{-t/(RC)}.
\]

El cambio ocurre cuando \(v_C=2V_{CC}/3\):

\[
\frac{2}{3}V_{CC}
=
V_{CC}-\frac{2}{3}V_{CC}e^{-t_H/(RC)}.
\]

\[
\frac{2}{3}=1-\frac{2}{3}e^{-t_H/(RC)}.
\]

\[
e^{-t_H/(RC)}=\frac{1}{2}.
\]

\[
t_H=RC\ln(2).
\]

Por simetria, al descargar entre \(2V_{CC}/3\) y \(V_{CC}/3\):

\[
t_L=RC\ln(2).
\]

Entonces:

\[
T=t_H+t_L=2RC\ln(2).
\]

\[
f_0=\frac{1}{2RC\ln(2)}.
\]

Se despeja:

\[
R_{osc}=\frac{1}{2f_0C_{osc}\ln(2)}.
\]

Con:

\[
C_{osc}=10~\text{nF}.
\]

\[
R_{osc}=
\frac{1}
{2(2290)(10\times10^{-9})\ln(2)}.
\]

\[
R_{osc}=31499.9~\Omega.
\]

\[
\boxed{R_{osc}\approx31.5~\text{k}\Omega}.
\]

En LTspice se usa \(R_{osc}=31.46~\text{k}\Omega\). La simulacion mide:

\[
f_{square}=2290.02~\text{Hz},
\qquad
D=50.00\%.
\]

\section{Onda diente de sierra positiva}

La diente de sierra se obtiene cargando un capacitor desde \(+15~\text{V}\) y
descargandolo a tierra en cada flanco de la onda cuadrada. Como se descarga en
ambos flancos:

Para detectar ambos flancos se usa la senal cuadrada directa y una version
invertida. La inversion se puede implementar con una etapa \texttt{LF353}; el
reset final lo realiza un transistor \texttt{2N3904}, por lo que la rampa no
depende de una fuente ideal de corriente.

\[
f_{saw}=2f_0=4580~\text{Hz}.
\]

\[
T_{saw}=\frac{1}{f_{saw}}
=
\frac{1}{4580}.
\]

\[
T_{saw}=2.1834\times10^{-4}~\text{s}.
\]

La carga del capacitor es:

\[
v_C(t)=V_{CC}\left(1-e^{-t/(RC)}\right).
\]

Se desea:

\[
v_C(T_{saw})=V_{saw,pk}.
\]

Por tanto:

\[
V_{saw,pk}=V_{CC}\left(1-e^{-T_{saw}/(RC)}\right).
\]

Se despeja \(R\):

\[
\frac{V_{saw,pk}}{V_{CC}}
=
1-e^{-T_{saw}/(RC)}.
\]

\[
e^{-T_{saw}/(RC)}
=
1-\frac{V_{saw,pk}}{V_{CC}}.
\]

\[
-\frac{T_{saw}}{RC}
=
\ln\left(1-\frac{V_{saw,pk}}{V_{CC}}\right).
\]

\[
R=
\frac{-T_{saw}}
{C\ln\left(1-\frac{V_{saw,pk}}{V_{CC}}\right)}.
\]

Con:

\[
V_{CC}=15~\text{V},
\qquad
V_{saw,pk}=5~\text{V},
\qquad
C_{saw}=10~\text{nF}.
\]

\[
R_{charge}
=
\frac{-2.1834\times10^{-4}}
{(10\times10^{-9})\ln\left(1-\frac{5}{15}\right)}.
\]

\[
R_{charge}=53849~\Omega.
\]

\[
\boxed{R_{charge}\approx53.8~\text{k}\Omega}.
\]

Se usa:

\[
\boxed{R_{charge}=53.6~\text{k}\Omega}.
\]

\section{Onda triangular sin flip-flop externo}

Para no usar un flip-flop externo, se usa un segundo \texttt{NE555} en modo
monoestable como divisor por 2. El monoestable se dispara con un flanco de la
cuadrada principal. Mientras esta temporizando, ignora el siguiente disparo.
Por eso vuelve a activarse cada dos periodos del oscilador.

El periodo del oscilador principal es:

\[
T_0=\frac{1}{f_0}.
\]

\[
T_0=\frac{1}{2290}=436.68~\mu\text{s}.
\]

Se escoge un tiempo monoestable un poco mayor que \(T_0\):

\[
t_w=1.05T_0.
\]

\[
t_w=1.05(436.68~\mu\text{s}).
\]

\[
t_w=458.52~\mu\text{s}.
\]

Como el segundo disparo llega cuando el monoestable aun esta activo, se ignora.
El siguiente disparo aceptado ocurre dos periodos despues:

\[
T_{div}=2T_0.
\]

\[
f_{div}=\frac{1}{T_{div}}=\frac{f_0}{2}.
\]

\[
\boxed{f_{div}=1145~\text{Hz}}.
\]

Para el \texttt{NE555} monoestable:

\[
t_w=1.1R_{div}C_{div}.
\]

Se despeja:

\[
R_{div}=\frac{t_w}{1.1C_{div}}.
\]

Con:

\[
C_{div}=10~\text{nF}.
\]

\[
R_{div}=
\frac{458.52\times10^{-6}}
{1.1(10\times10^{-9})}.
\]

\[
R_{div}=41684~\Omega.
\]

\[
\boxed{R_{div}\approx41.7~\text{k}\Omega}.
\]

El ciclo alto del divisor es:

\[
D=\frac{t_w}{2T_0}.
\]

\[
D=\frac{1.05T_0}{2T_0}=0.525.
\]

Como el ciclo no es exactamente \(50\%\), el cambiador de nivel debe producir
una senal de promedio cero antes del integrador. Se escoge:

\[
V_{high}=+1.00~\text{V}.
\]

La condicion de promedio cero es:

\[
DV_{high}+(1-D)V_{low}=0.
\]

Se despeja:

\[
V_{low}=-\frac{D}{1-D}V_{high}.
\]

\[
V_{low}=-\frac{0.525}{0.475}(1.00).
\]

\[
\boxed{V_{low}=-1.105~\text{V}}.
\]

Por tanto, el cambiador de nivel convierte:

\[
0~\text{V}\rightarrow -1.105~\text{V},
\qquad
15~\text{V}\rightarrow +1.000~\text{V}.
\]

Este cambiador de nivel se implementa con \texttt{LF353}. La misma familia de
amplificador operacional se usa luego para el integrador inversor que genera la
senal triangular.

Una expresion lineal para este bloque es:

\[
V_{drive}=kV_{div}-V_{ref}.
\]

\[
k=\frac{1.000-(-1.105)}{15}=0.1403.
\]

\[
V_{ref}=1.105~\text{V}.
\]

Para el integrador inversor:

\[
\Delta V=\frac{V_{high}t_w}{R_{tri}C_{tri}}.
\]

Se desea que el cambio de minimo a maximo sea:

\[
\Delta V=2V_{tri,pk}.
\]

Entonces:

\[
R_{tri}=
\frac{V_{high}t_w}
{2V_{tri,pk}C_{tri}}.
\]

Con:

\[
V_{tri,pk}=2~\text{V},
\qquad
C_{tri}=10~\text{nF}.
\]

\[
R_{tri}=
\frac{(1.00)(458.52\times10^{-6})}
{2(2)(10\times10^{-9})}.
\]

\[
R_{tri}=11463~\Omega.
\]

\[
\boxed{R_{tri}\approx11.5~\text{k}\Omega}.
\]

La simulacion del divisor con el segundo \texttt{NE555} mide:

\[
f_{div}=1145.00~\text{Hz}.
\]

\section{Onda senoidal por filtro pasabajas}

\subsection{Seleccion de la entrada}

El filtro pasabajas debe dejar pasar principalmente la frecuencia de interes y
atenuar las variaciones mas rapidas. En terminos de Bode, se busca mantener la
senal en la banda de paso y rechazar con fuerza las componentes ubicadas por
encima de la frecuencia de corte.

Entre las tres senales previas:

\[
\begin{array}{c|c}
\text{Senal} & \text{Caracteristica temporal} \\
\hline
\text{Cuadrada} & \text{cambia de nivel bruscamente dos veces por periodo}\\
\text{Diente de sierra} & \text{sube suavemente, pero regresa con un salto}\\
\text{Triangular} & \text{no tiene salto de voltaje; solo cambia de pendiente}
\end{array}
\]

Una senoide tampoco tiene saltos de voltaje. Por tanto, la triangular es la
entrada mas conveniente para el pasabajas, porque antes del filtro ya es la
senal mas continua de las tres.

\subsection{Tipo de respuesta}

Segun el enunciado:

\[
\begin{array}{c|c}
Z & \text{Respuesta} \\
\hline
1,2,3 & \text{Bessel}\\
4,5,6 & \text{Butterworth}\\
7,8,9 & \text{Chebyshev}
\end{array}
\]

Como:

\[
Z=9,
\]

se usa:

\[
\boxed{\text{Chebyshev tipo I}}.
\]

La materia indica que Chebyshev tiene una transicion mas abrupta entre banda de
paso y banda de rechazo, a cambio de rizado en la banda de paso.

\subsection{Frecuencia de corte}

La frecuencia de la triangular es:

\[
f_{tri}=1145~\text{Hz}.
\]

Se escoge:

\[
f_c=1.25f_{tri}.
\]

\[
f_c=1.25(1145).
\]

\[
\boxed{f_c=1431.25~\text{Hz}}.
\]

Asi, \(f_{tri}\) queda dentro de la banda de paso.

\subsection{Estructura del filtro}

La materia de filtros de orden superior usa etapas en cascada. Para un filtro
de quinto orden:

Las etapas activas del filtro se implementan con \texttt{LF353}. Se usa una
seccion de primer orden y dos secciones Sallen-Key de segundo orden.

\[
5=1+2+2.
\]

Por tanto, se implementa:

\[
\text{primer orden} + \text{segundo orden} + \text{segundo orden}.
\]

La frecuencia normalizada es:

\[
S=\frac{j\omega}{\omega_0}=\frac{jf}{f_0}.
\]

Para una etapa de segundo orden:

\[
H(S)=H_0\frac{N(S)}{1+aS+bS^2}.
\]

Para pasabajas:

\[
N(S)=1.
\]

El factor de calidad se calcula como:

\[
Q=\frac{\sqrt{b}}{a}.
\]

\subsection{Coeficientes tabulados}

Para Chebyshev tipo I de quinto orden se usan tres etapas:

\[
\begin{array}{c|c|c|c}
\text{Etapa} & \text{Orden} & a_i & b_i \\
\hline
1 & 1 & 2.7600 & 0\\
2 & 2 & 1.2296 & 2.0975\\
3 & 2 & 0.2162 & 0.9655
\end{array}
\]

\subsection{Etapa de primer orden}

La primera etapa tiene:

\[
H_1(S)=H_{01}\frac{1}{1+a_1S}.
\]

La frecuencia del polo fisico es:

\[
f_{p1}=\frac{f_c}{a_1}.
\]

\[
f_{p1}=\frac{1431.25}{2.7600}.
\]

\[
f_{p1}=518.57~\text{Hz}.
\]

Con:

\[
C_f=10~\text{nF},
\]

se obtiene:

\[
R_{p1}=\frac{1}{2\pi f_{p1}C_f}.
\]

\[
R_{p1}=
\frac{1}{2\pi(518.57)(10\times10^{-9})}.
\]

\[
R_{p1}=30.7~\text{k}\Omega.
\]

\subsection{Primera etapa de segundo orden}

De la tabla:

\[
a_2=1.2296,
\qquad
b_2=2.0975.
\]

El factor de calidad es:

\[
Q_2=\frac{\sqrt{b_2}}{a_2}.
\]

\[
Q_2=\frac{\sqrt{2.0975}}{1.2296}.
\]

\[
Q_2=1.178.
\]

La frecuencia natural es:

\[
f_{n2}=\frac{f_c}{\sqrt{b_2}}.
\]

\[
f_{n2}=\frac{1431.25}{\sqrt{2.0975}}.
\]

\[
f_{n2}=988.25~\text{Hz}.
\]

La resistencia para Sallen-Key con componentes iguales es:

\[
R_{SK1}=\frac{1}{2\pi f_{n2}C_f}.
\]

\[
R_{SK1}=
\frac{1}{2\pi(988.25)(10\times10^{-9})}.
\]

\[
R_{SK1}=16.1~\text{k}\Omega.
\]

Para Sallen-Key con \(R_1=R_2\) y \(C_1=C_2\):

\[
Q=\frac{1}{3-K}.
\]

Entonces:

\[
K=3-\frac{1}{Q}.
\]

\[
K_1=3-\frac{1}{1.178}.
\]

\[
K_1=2.151.
\]

Con:

\[
R_g=10~\text{k}\Omega,
\]

\[
R_{f1}=(K_1-1)R_g.
\]

\[
R_{f1}=(2.151-1)(10~\text{k}\Omega).
\]

\[
R_{f1}=11.5~\text{k}\Omega.
\]

\subsection{Segunda etapa de segundo orden}

De la tabla:

\[
a_3=0.2162,
\qquad
b_3=0.9655.
\]

\[
Q_3=\frac{\sqrt{b_3}}{a_3}.
\]

\[
Q_3=\frac{\sqrt{0.9655}}{0.2162}.
\]

\[
Q_3=4.545.
\]

\[
f_{n3}=\frac{f_c}{\sqrt{b_3}}.
\]

\[
f_{n3}=\frac{1431.25}{\sqrt{0.9655}}.
\]

\[
f_{n3}=1456.63~\text{Hz}.
\]

\[
R_{SK2}=\frac{1}{2\pi f_{n3}C_f}.
\]

\[
R_{SK2}=
\frac{1}{2\pi(1456.63)(10\times10^{-9})}.
\]

\[
R_{SK2}=10.9~\text{k}\Omega.
\]

\[
K_2=3-\frac{1}{Q_3}.
\]

\[
K_2=3-\frac{1}{4.545}.
\]

\[
K_2=2.780.
\]

\[
R_{f2}=(K_2-1)R_g.
\]

\[
R_{f2}=(2.780-1)(10~\text{k}\Omega).
\]

\[
R_{f2}=17.8~\text{k}\Omega.
\]

\subsection{Atenuador de entrada}

La ganancia de las dos etapas Sallen-Key es:

\[
K_T=K_1K_2.
\]

\[
K_T=(2.151)(2.780).
\]

\[
K_T=5.98.
\]

Para compensarla:

\[
A=\frac{1}{K_T}.
\]

Usando un divisor de entrada:

\[
A=\frac{R_{bot}}{R_{top}+R_{bot}}.
\]

Se despeja:

\[
R_{top}=R_{bot}(K_T-1).
\]

Con:

\[
R_{bot}=100~\text{k}\Omega,
\]

\[
R_{top}=(100~\text{k}\Omega)(5.98-1).
\]

\[
R_{top}=498~\text{k}\Omega.
\]

\subsection{Valores del filtro}

Con \(C_f=10~\text{nF}\), los valores usados son:

\[
\begin{array}{c|c}
\text{Bloque} & \text{Valores} \\
\hline
\text{Atenuador de entrada} & R_{top}=498~\text{k}\Omega,\ R_{bot}=100~\text{k}\Omega\\
\text{Primer orden} & R=30.7~\text{k}\Omega,\ C=10~\text{nF}\\
\text{Sallen-Key 1} & R_1=R_2=16.1~\text{k}\Omega,\ C_1=C_2=10~\text{nF}\\
\text{Ganancia SK1} & R_g=10~\text{k}\Omega,\ R_f=11.5~\text{k}\Omega\\
\text{Sallen-Key 2} & R_1=R_2=10.9~\text{k}\Omega,\ C_1=C_2=10~\text{nF}\\
\text{Ganancia SK2} & R_g=10~\text{k}\Omega,\ R_f=17.8~\text{k}\Omega
\end{array}
\]

\subsection{Verificacion por Bode}

La simulacion AC del filtro da:

\[
|H(1145~\text{Hz})|=-11.12~\text{dB}.
\]

\[
|H(3435~\text{Hz})|=-62.29~\text{dB}.
\]

La diferencia entre ambas ganancias es:

\[
62.29-11.12=51.17~\text{dB}.
\]

Por tanto, el filtro deja pasar mucho mas la frecuencia de interes que una
frecuencia tres veces mayor. Esto valida el comportamiento pasabajas usando
diagramas de Bode.

\section{Resumen numerico}

\[
\begin{array}{c|c}
\text{Parametro} & \text{Valor} \\
\hline
f_0 & 2290~\text{Hz}\\
f_{saw} & 4580~\text{Hz}\\
f_{tri} & 1145~\text{Hz}\\
R_{osc} & 31.46~\text{k}\Omega\\
R_{charge} & 53.6~\text{k}\Omega\\
R_{div} & 41.7~\text{k}\Omega\\
R_{tri} & 11.5~\text{k}\Omega\\
f_c & 1431.25~\text{Hz}\\
Opamp usado & \text{LF353}\\
R_{p1} & 30.7~\text{k}\Omega\\
R_{SK1} & 16.1~\text{k}\Omega\\
R_{f1} & 11.5~\text{k}\Omega\\
R_{SK2} & 10.9~\text{k}\Omega\\
R_{f2} & 17.8~\text{k}\Omega
\end{array}
\]

\end{document}
